\section{Discussion}

The prescribed-fiber theorem made the first Hasse example almost automatic:
start with an intersective polynomial and compile its finite \'etale algebra
into a full Keller fiber.  That flexibility also concealed the next
question.  If the construction is allowed to absorb the polynomial into the
map, it says nothing about how many arithmetic behaviors one map can carry.

The present argument began when those two roles were separated.  The map
should remember only the fixed quintic shape; the target should remember the
parameter.  The quantifier change is small to write and substantial in
effect:
\[
\begin{array}{c|c|c}
 &\text{map}&\text{fiber algebra}\\ \hline
\text{prescribed fibers}&\text{varies}&\text{arbitrary of rank }\ge3\\
\text{this paper}&\text{fixed}&
 \Q[X]/((X^3-a)(X^2+X+1)),\ a\in\mathcal A.
\end{array}
\tag{7.1}
\]

Two observations make the fixed row possible.  The first is geometric:
after the translation \(X=S-\tfrac12\), the quadratic
\(X^2+X+1\) becomes \(S^2+\tfrac34\).  The parameter \(a\) then moves
only the quadratic and constant coefficients, which are exactly the two
free coefficients of the quadratic-gauge target pencil.  The second is
arithmetic: once the local proof is written place by place, primality is
nowhere essential.  The congruence-and-support core is multiplicatively
closed; removing its perfect cubes leaves the admissible family.  The
count grows from the prime scale \(B/\log B\) to
\(B/\sqrt{\log B}\).

There is also a useful warning in the construction.  The target line is
exceptional from the viewpoint of generic monodromy.  Every inverse
polynomial on it has the same quadratic-times-cubic factorization.  This is
not an accidental defect: the factorization is what lets one component
supply local points precisely where the other does not.  The deterministic
local behavior on the line therefore says little about the distribution of
local points on the ambient three-dimensional target.

The exact scope is equally important.  Theorem~\ref{thm:count-intro} counts
one sufficient multiplicative family.  It does not classify all Hasse
failures on the line, and still less all Hasse-failing fibers of the map.
A converse local analysis would have to decide which arbitrary rational
parameters make at least one of the two factors soluble at every place.
Integral source points introduce another layer: a local root of the inverse
polynomial need not reconstruct to an integral source point without further
valuation control.

What the paper does establish is a clean progression.  What began with one
stubborn quintic ends with a more generous picture: the classical example is
the first visible point on a rational target line belonging to one
Jacobian-one map of the smallest possible geometric degree.  A half-unit
translation reveals the line, a place-by-place proof reveals the
multiplicative family, and the arithmetic itself explains why degree five is
minimal.  These are small observations individually, but together they turn
an example into a landscape.

\paragraph{Provenance and AI assistance.}
The arithmetic polynomial at \(a=19\) is due to Berend and Bilu
\cite{BerendBilu}.  The fixed-map construction is obtained by specializing
the quadratic-gauge prescribed-fiber mechanism of \cite{vanRijnFibers}.

I used generative AI tools during exploration, symbolic-checker development,
proof review, literature search, and manuscript editing.  I checked the
stated mathematical claims and take responsibility for the paper.
